\section{Sistemas de unidades}

El objetivo de un sistema de unidades es el de definir en términos de una unidad estandarizada la magnitud de una cantidad medida.

En 1960, el Comité Internacional de Pesos y Medidas proporciona definiciones claras respecto a unidades estándar, ayudando a los científicos para que comuniquen sus medidas de manera más precisa. El Sistema Internacional (SI) es el más adoptado en el mundo.

\subsection{Magnitudes dundamentales y derivadas}



    \begin{table}[t]
    \centering
    \begin{threeparttable}
    \caption{Magnitudes fundamentales}
    \label{magn_fund}
    {\setstretch{1.25}
        \begin{tabular}{|p{0.25\textwidth}|p{0.175\textwidth}|p{0.175\textwidth}|} \hline
            \bfseries{Dimensión} & \textbf{Símbolo} & \textbf{Unidad} \\ \hline
            Longitud & L & Metro ($m$) \\ \hline
            Masa & M & Kilogramo ($kg$) \\ \hline
            Tiempo & T & Segundo ($s$) \\ \hline
            Intensidad & I & A \\ \hline
            Temperatuda & θ & Kelvin ($K$) \\ \hline
            Intensidad lumínica & J & Candela ($cd$) \\ \hline
            Cantidad de sustancia & N & Mol \\ \hline
        \end{tabular}}
        \begin{tablenotes} % \tnote{#}
            \item[] 
        \end{tablenotes}
    \end{threeparttable}
\end{table}
    \begin{table}[t]
    \centering
    \begin{threeparttable}
    \caption{Magnitudes derivadas}
    \label{magn_der}
    {\setstretch{1.25}
        \begin{tabular}{|p{0.25\textwidth}|p{0.175\textwidth}|p{0.175\textwidth}|} \hline
            \bfseries{Dimensión} & \textbf{Símbolo} & \textbf{Unidad} \\ \hline
            Área & $L^2$ & $m^2$ \\ \hline
            Volumen & $L^3$ & $m^3$ \\ \hline
            Velocidad & $\frac{L}{T}$ & $\frac{m}{s}$ \\ \hline
            Aceleración & $\frac{L}{T^2}$ & $\frac{m}{s^2}$ \\ \hline
            Fuerza & $\frac{M*L}{T^2}$ & $\frac{kg*m}{s^{2}}=N$ \\ \hline
            Momento & $\frac{M*L^2}{T^2}$ & $n*m=J$ \\ \hline
            Trabajo & $\frac{M*L^2}{T^2}$ & $J$ \\ \hline
            Potencia & $\frac{M*L^2}{T^3}$ & $\frac{J}{s}$ \\ \hline
            Densidad & $\frac{M}{L^3}$ & $\frac{kg}{m^3}$ \\ \hline
            Peso & $\frac{M}{L^2*T^2}$ & $\frac{N}{m^3}$ \\ \hline
            Impulso & $\frac{M*L}{T}$ & $N*S$ \\ \hline
        \end{tabular}}
        \begin{tablenotes} % \tnote{#}
            \item[] 
        \end{tablenotes}
    \end{threeparttable}
\end{table}

\subsection{Unidades básicas y prácticas}



Unidades básicas o patrón:

\begin{itemize}
    \item Masa ($kg$)
    \item Longitud ($m$)
    \item Tiempo ($s$).
\end{itemize}

Unidades prácticas:

\begin{itemize}
    \item Masa ($g$)
    \item Longitud ($km$)
    \item Tiempo ($h$).
\end{itemize}

\section{Notación científica}

La notación científica, también llamada notación exponencial, es una forma de escribir números basada en potencias de 10. Esta se representa con la forma $m \times 10^{n}$, donde:

\begin{itemize}
    \item $m$, llamada mantisa, es un número decimal cuyo entero es una sola cifra distinta de 0
    \item $10^{n}$ es la potencia u orden de magnitud, y es un exponente entero que indica 
\end{itemize}

Esta forma de representar números surgida en el siglo III a.C. y descrita por Arquímedes en su obra El Contador de Arena, es muy conveniente para escribir números muy grandes o pequeños y hacer cálculos con ellos, lo cual es muy útil en materias donde se realizan mediciones muy grandes o muy pequeñas, como astronomía o microbiología.

Algunos ejemplos de notación científica son los siguientes:

\begin{enumerate}
    \item Masa del protón: $1.67 \times 10^{-27}$
    \item Masa del electrón: $9.11 \times 10^{-31}$
    \item Número de Avogadro: $6.02 \times 10^{23}$
    \item 1 año luz: $1 \times 10^{16}m$
    \item Diámetro del leucocito: $6.5 \times 10^{-3}cm$
\end{enumerate}

\section{}



\section{}

